% =========================================================
% ppt/src/minggu-12.tex
% Slide pertemuan 12
% Bab XIII — Prosedur dan Manajemen Stack
%
% Sub-CPMK 3.4 | CPMK03 | CPL03
% =========================================================

\documentclass[aspectratio=169,11pt]{beamer}
% =========================================================
% ppt/src/preamble-ppt.tex
% Tema, warna, paket, dan makro bersama seluruh deck slide.
%
% Dipakai oleh setiap berkas minggu-NN.tex dengan pola:
%   \documentclass[aspectratio=169,11pt]{beamer}
%   % =========================================================
% ppt/src/preamble-ppt.tex
% Tema, warna, paket, dan makro bersama seluruh deck slide.
%
% Dipakai oleh setiap berkas minggu-NN.tex dengan pola:
%   \documentclass[aspectratio=169,11pt]{beamer}
%   % =========================================================
% ppt/src/preamble-ppt.tex
% Tema, warna, paket, dan makro bersama seluruh deck slide.
%
% Dipakai oleh setiap berkas minggu-NN.tex dengan pola:
%   \documentclass[aspectratio=169,11pt]{beamer}
%   \input{preamble-ppt}
%
% Palet warna disalin dari buku_ajar/config/styles.tex agar
% identitas visual slide sama dengan buku ajar.
% =========================================================

% ---------------------------------------------------------
% 1. Paket
% ---------------------------------------------------------
\usepackage[T1]{fontenc}
\usepackage[utf8]{inputenc}
\usepackage[indonesian]{babel}
\usepackage{lmodern}
\usepackage{tikz}
\usetikzlibrary{arrows.meta, positioning, calc, shapes.geometric}
\usepackage{booktabs}
\usepackage{array}
\usepackage{amsmath}
\usepackage{amsfonts}
\usepackage{pifont}          % \ding untuk tanda checklist
\usepackage{listings}

% ---------------------------------------------------------
% 2. Palet warna OBE (identik dengan buku)
% ---------------------------------------------------------
\definecolor{obeBlue}{HTML}{1F3864}
\definecolor{obeBlueLight}{HTML}{DCE4F5}
\definecolor{obeTeal}{HTML}{0F6E6E}
\definecolor{obeTealLight}{HTML}{D9F0EE}
\definecolor{obeGreen}{HTML}{2E6B34}
\definecolor{obeGreenLight}{HTML}{DEEFDF}
\definecolor{obeAmber}{HTML}{8A5A00}
\definecolor{obeAmberLight}{HTML}{FDF1D6}
\definecolor{obeMaroon}{HTML}{7A2233}
\definecolor{obeMaroonLight}{HTML}{F7E1E5}
\definecolor{obeGray}{HTML}{4A4A4A}
\definecolor{obeGrayLight}{HTML}{F0F0F0}
\definecolor{codeBack}{HTML}{F7F7F4}
\definecolor{codeFrame}{HTML}{C9C9C0}

% ---------------------------------------------------------
% 3. Identitas mata kuliah (dipakai seluruh deck)
% ---------------------------------------------------------
\newcommand{\NamaMK}{Bahasa Pemrograman Tingkat Rendah}
\newcommand{\NamaMKShort}{Bhs.\ Pemrograman Tingkat Rendah}
\newcommand{\KodeMK}{MWCA010}
\newcommand{\Prodi}{Program Studi Teknik Informatika}
\newcommand{\Fakultas}{Fakultas Teknologi Informasi}
\newcommand{\PerguruanTinggi}{Universitas Bale Bandung}
\newcommand{\DosenPengampu}{Cecep Suwanda, S.Si., M.Kom.}
\newcommand{\SemesterTeks}{Semester Ganjil}

% \PertemuanTeks dan \JudulBab ditetapkan di masing-masing
% berkas minggu-NN.tex sebelum \begin{document}.
\providecommand{\PertemuanTeks}{Pertemuan}

% ---------------------------------------------------------
% 4. Tema beamer
% ---------------------------------------------------------
\usetheme{default}
\setbeamertemplate{navigation symbols}{}

\setbeamercolor{normal text}{fg=black!88}
\setbeamercolor{structure}{fg=obeBlue}
\setbeamercolor{alerted text}{fg=obeMaroon}

% --- judul slide: pita biru penuh lebar ---
\setbeamercolor{frametitle}{fg=white,bg=obeBlue}
\setbeamerfont{frametitle}{series=\bfseries,size=\large}
\setbeamertemplate{frametitle}{%
    \nointerlineskip%
    \begin{beamercolorbox}[%
        wd=\paperwidth,ht=2.9ex,dp=1.5ex,
        leftskip=0.4cm,rightskip=0.4cm]{frametitle}%
        \usebeamerfont{frametitle}\insertframetitle%
    \end{beamercolorbox}%
    \nointerlineskip%
    \hbox{\textcolor{obeTeal}{\rule{\paperwidth}{1.1pt}}}%
}

% --- kaki slide: identitas mata kuliah, pertemuan, nomor ---
\setbeamercolor{footline}{fg=white,bg=obeGray}
\setbeamerfont{footline}{size=\scriptsize,family=\sffamily}
\setbeamertemplate{footline}{%
    \leavevmode%
    \hbox{%
        \begin{beamercolorbox}[wd=0.42\paperwidth,ht=2.1ex,dp=1.1ex,
            leftskip=0.4cm]{footline}%
            \usebeamerfont{footline}%
            \KodeMK\ \textbullet\ \NamaMKShort%
        \end{beamercolorbox}%
        \begin{beamercolorbox}[wd=0.32\paperwidth,ht=2.1ex,dp=1.1ex,
            center]{footline}%
            \usebeamerfont{footline}\PertemuanTeks%
        \end{beamercolorbox}%
        \begin{beamercolorbox}[wd=0.26\paperwidth,ht=2.1ex,dp=1.1ex,
            rightskip=0.4cm]{footline}%
            \usebeamerfont{footline}%
            \insertframenumber\,/\,\inserttotalframenumber%
        \end{beamercolorbox}%
    }%
    \vskip0pt%
}

% --- daftar ---
\setbeamercolor{itemize item}{fg=obeTeal}
\setbeamercolor{itemize subitem}{fg=obeTeal}
\setbeamercolor{enumerate item}{fg=obeBlue}
\setbeamertemplate{itemize item}{\raise0.2ex\hbox{\tiny$\blacksquare$}}
\setbeamertemplate{itemize subitem}{\raise0.2ex\hbox{\scriptsize$\bullet$}}

% --- kotak (pengganti tcolorbox pada buku) ---
\setbeamercolor{block title}{bg=obeBlueLight,fg=obeBlue}
\setbeamercolor{block body}{bg=obeBlueLight!45,fg=black!88}
\setbeamercolor{block title alerted}{bg=obeMaroonLight,fg=obeMaroon}
\setbeamercolor{block body alerted}{bg=obeMaroonLight!45,fg=black!88}
\setbeamercolor{block title example}{bg=obeTealLight,fg=obeTeal}
\setbeamercolor{block body example}{bg=obeTealLight!45,fg=black!88}

% ---------------------------------------------------------
% 5. Makro notasi (sama persis dengan buku)
%
% Catatan: \note TIDAK didefinisikan di sini karena beamer
% sudah memakai \note untuk catatan pembicara. Untuk catatan
% di slide dipakai \catatan.
% ---------------------------------------------------------
\newcommand{\code}[1]{\texttt{#1}}
\newcommand{\asm}[1]{\texttt{#1}}
\newcommand{\reg}[1]{\texttt{#1}}
\newcommand{\flag}[1]{\texttt{#1}}
\newcommand{\dir}[1]{\texttt{#1}}
\newcommand{\file}[1]{\texttt{#1}}
\newcommand{\eng}[1]{\textit{#1}}

\newcommand{\hex}[1]{\texttt{#1h}}
\newcommand{\bin}[1]{\texttt{#1b}}
\newcommand{\dec}[1]{\texttt{#1}}

\newcommand{\important}[1]{\textbf{Penting:} #1}
\newcommand{\catatan}[1]{\textbf{Catatan:} #1}
\newcommand{\term}[1]{\textbf{#1}}

\newcommand{\CPL}[1]{\textbf{CPL#1}}
\newcommand{\CPMK}[1]{\textbf{CPMK#1}}
\newcommand{\SubCPMK}[1]{\textbf{Sub-CPMK #1}}

\newcommand{\checkbox}{\ding{113}\hspace{0.35em}}
\newcommand{\checked}{\ding{51}\hspace{0.35em}}
\newcommand{\crossed}{\ding{55}\hspace{0.35em}}

\newcommand{\flowarrow}{\ensuremath{\rightarrow}}

% ---------------------------------------------------------
% 6. Makro tata letak slide
% ---------------------------------------------------------

% Keterangan sumber materi (ketertelusuran ke buku)
\newcommand{\Sumber}[1]{%
    \vfill
    {\scriptsize\color{obeGray}Sumber: #1\par}%
}

% Frame judul deck: \JudulDeck{<judul bab>}{<pertemuan>}
\newcommand{\JudulDeck}[2]{%
\begin{frame}[plain]
    \vspace{0.25cm}
    {\color{obeBlue}\sffamily\bfseries\small \NamaMK\ (\KodeMK)\par}
    \vspace{0.12cm}
    {\color{obeGray}\sffamily\scriptsize
     \Prodi\ \textbullet\ \Fakultas\ \textbullet\ \PerguruanTinggi\par}
    \vfill
    {\color{obeBlue}\sffamily\LARGE\bfseries #1\par}
    \vspace{0.4cm}
    {\color{obeTeal}\sffamily\large #2\par}
    \vspace{0.45cm}
    {\color{obeTeal}\rule{\textwidth}{1pt}}
    \vfill
    {\color{obeGray}\sffamily\small Dosen Pengampu\par}
    \vspace{0.08cm}
    {\sffamily\bfseries \DosenPengampu\par}
    \vspace{0.2cm}
    {\color{obeGray}\sffamily\small \SemesterTeks\par}
    \vspace{0.25cm}
\end{frame}
}

% Frame pemisah bagian
\newcommand{\PemisahBagian}[1]{%
\begin{frame}[plain]
    \vfill
    \begin{center}
        {\color{obeTeal}\sffamily\small\bfseries BAGIAN\par}
        \vspace{0.35cm}
        {\color{obeBlue}\sffamily\Large\bfseries #1\par}
        \vspace{0.35cm}
        {\color{obeTeal}\rule{0.55\textwidth}{0.9pt}}
    \end{center}
    \vfill
\end{frame}
}

% ---------------------------------------------------------
% 7. Gaya penampil kode
% ---------------------------------------------------------
\input{listing-ppt}

%
% Palet warna disalin dari buku_ajar/config/styles.tex agar
% identitas visual slide sama dengan buku ajar.
% =========================================================

% ---------------------------------------------------------
% 1. Paket
% ---------------------------------------------------------
\usepackage[T1]{fontenc}
\usepackage[utf8]{inputenc}
\usepackage[indonesian]{babel}
\usepackage{lmodern}
\usepackage{tikz}
\usetikzlibrary{arrows.meta, positioning, calc, shapes.geometric}
\usepackage{booktabs}
\usepackage{array}
\usepackage{amsmath}
\usepackage{amsfonts}
\usepackage{pifont}          % \ding untuk tanda checklist
\usepackage{listings}

% ---------------------------------------------------------
% 2. Palet warna OBE (identik dengan buku)
% ---------------------------------------------------------
\definecolor{obeBlue}{HTML}{1F3864}
\definecolor{obeBlueLight}{HTML}{DCE4F5}
\definecolor{obeTeal}{HTML}{0F6E6E}
\definecolor{obeTealLight}{HTML}{D9F0EE}
\definecolor{obeGreen}{HTML}{2E6B34}
\definecolor{obeGreenLight}{HTML}{DEEFDF}
\definecolor{obeAmber}{HTML}{8A5A00}
\definecolor{obeAmberLight}{HTML}{FDF1D6}
\definecolor{obeMaroon}{HTML}{7A2233}
\definecolor{obeMaroonLight}{HTML}{F7E1E5}
\definecolor{obeGray}{HTML}{4A4A4A}
\definecolor{obeGrayLight}{HTML}{F0F0F0}
\definecolor{codeBack}{HTML}{F7F7F4}
\definecolor{codeFrame}{HTML}{C9C9C0}

% ---------------------------------------------------------
% 3. Identitas mata kuliah (dipakai seluruh deck)
% ---------------------------------------------------------
\newcommand{\NamaMK}{Bahasa Pemrograman Tingkat Rendah}
\newcommand{\NamaMKShort}{Bhs.\ Pemrograman Tingkat Rendah}
\newcommand{\KodeMK}{MWCA010}
\newcommand{\Prodi}{Program Studi Teknik Informatika}
\newcommand{\Fakultas}{Fakultas Teknologi Informasi}
\newcommand{\PerguruanTinggi}{Universitas Bale Bandung}
\newcommand{\DosenPengampu}{Cecep Suwanda, S.Si., M.Kom.}
\newcommand{\SemesterTeks}{Semester Ganjil}

% \PertemuanTeks dan \JudulBab ditetapkan di masing-masing
% berkas minggu-NN.tex sebelum \begin{document}.
\providecommand{\PertemuanTeks}{Pertemuan}

% ---------------------------------------------------------
% 4. Tema beamer
% ---------------------------------------------------------
\usetheme{default}
\setbeamertemplate{navigation symbols}{}

\setbeamercolor{normal text}{fg=black!88}
\setbeamercolor{structure}{fg=obeBlue}
\setbeamercolor{alerted text}{fg=obeMaroon}

% --- judul slide: pita biru penuh lebar ---
\setbeamercolor{frametitle}{fg=white,bg=obeBlue}
\setbeamerfont{frametitle}{series=\bfseries,size=\large}
\setbeamertemplate{frametitle}{%
    \nointerlineskip%
    \begin{beamercolorbox}[%
        wd=\paperwidth,ht=2.9ex,dp=1.5ex,
        leftskip=0.4cm,rightskip=0.4cm]{frametitle}%
        \usebeamerfont{frametitle}\insertframetitle%
    \end{beamercolorbox}%
    \nointerlineskip%
    \hbox{\textcolor{obeTeal}{\rule{\paperwidth}{1.1pt}}}%
}

% --- kaki slide: identitas mata kuliah, pertemuan, nomor ---
\setbeamercolor{footline}{fg=white,bg=obeGray}
\setbeamerfont{footline}{size=\scriptsize,family=\sffamily}
\setbeamertemplate{footline}{%
    \leavevmode%
    \hbox{%
        \begin{beamercolorbox}[wd=0.42\paperwidth,ht=2.1ex,dp=1.1ex,
            leftskip=0.4cm]{footline}%
            \usebeamerfont{footline}%
            \KodeMK\ \textbullet\ \NamaMKShort%
        \end{beamercolorbox}%
        \begin{beamercolorbox}[wd=0.32\paperwidth,ht=2.1ex,dp=1.1ex,
            center]{footline}%
            \usebeamerfont{footline}\PertemuanTeks%
        \end{beamercolorbox}%
        \begin{beamercolorbox}[wd=0.26\paperwidth,ht=2.1ex,dp=1.1ex,
            rightskip=0.4cm]{footline}%
            \usebeamerfont{footline}%
            \insertframenumber\,/\,\inserttotalframenumber%
        \end{beamercolorbox}%
    }%
    \vskip0pt%
}

% --- daftar ---
\setbeamercolor{itemize item}{fg=obeTeal}
\setbeamercolor{itemize subitem}{fg=obeTeal}
\setbeamercolor{enumerate item}{fg=obeBlue}
\setbeamertemplate{itemize item}{\raise0.2ex\hbox{\tiny$\blacksquare$}}
\setbeamertemplate{itemize subitem}{\raise0.2ex\hbox{\scriptsize$\bullet$}}

% --- kotak (pengganti tcolorbox pada buku) ---
\setbeamercolor{block title}{bg=obeBlueLight,fg=obeBlue}
\setbeamercolor{block body}{bg=obeBlueLight!45,fg=black!88}
\setbeamercolor{block title alerted}{bg=obeMaroonLight,fg=obeMaroon}
\setbeamercolor{block body alerted}{bg=obeMaroonLight!45,fg=black!88}
\setbeamercolor{block title example}{bg=obeTealLight,fg=obeTeal}
\setbeamercolor{block body example}{bg=obeTealLight!45,fg=black!88}

% ---------------------------------------------------------
% 5. Makro notasi (sama persis dengan buku)
%
% Catatan: \note TIDAK didefinisikan di sini karena beamer
% sudah memakai \note untuk catatan pembicara. Untuk catatan
% di slide dipakai \catatan.
% ---------------------------------------------------------
\newcommand{\code}[1]{\texttt{#1}}
\newcommand{\asm}[1]{\texttt{#1}}
\newcommand{\reg}[1]{\texttt{#1}}
\newcommand{\flag}[1]{\texttt{#1}}
\newcommand{\dir}[1]{\texttt{#1}}
\newcommand{\file}[1]{\texttt{#1}}
\newcommand{\eng}[1]{\textit{#1}}

\newcommand{\hex}[1]{\texttt{#1h}}
\newcommand{\bin}[1]{\texttt{#1b}}
\newcommand{\dec}[1]{\texttt{#1}}

\newcommand{\important}[1]{\textbf{Penting:} #1}
\newcommand{\catatan}[1]{\textbf{Catatan:} #1}
\newcommand{\term}[1]{\textbf{#1}}

\newcommand{\CPL}[1]{\textbf{CPL#1}}
\newcommand{\CPMK}[1]{\textbf{CPMK#1}}
\newcommand{\SubCPMK}[1]{\textbf{Sub-CPMK #1}}

\newcommand{\checkbox}{\ding{113}\hspace{0.35em}}
\newcommand{\checked}{\ding{51}\hspace{0.35em}}
\newcommand{\crossed}{\ding{55}\hspace{0.35em}}

\newcommand{\flowarrow}{\ensuremath{\rightarrow}}

% ---------------------------------------------------------
% 6. Makro tata letak slide
% ---------------------------------------------------------

% Keterangan sumber materi (ketertelusuran ke buku)
\newcommand{\Sumber}[1]{%
    \vfill
    {\scriptsize\color{obeGray}Sumber: #1\par}%
}

% Frame judul deck: \JudulDeck{<judul bab>}{<pertemuan>}
\newcommand{\JudulDeck}[2]{%
\begin{frame}[plain]
    \vspace{0.25cm}
    {\color{obeBlue}\sffamily\bfseries\small \NamaMK\ (\KodeMK)\par}
    \vspace{0.12cm}
    {\color{obeGray}\sffamily\scriptsize
     \Prodi\ \textbullet\ \Fakultas\ \textbullet\ \PerguruanTinggi\par}
    \vfill
    {\color{obeBlue}\sffamily\LARGE\bfseries #1\par}
    \vspace{0.4cm}
    {\color{obeTeal}\sffamily\large #2\par}
    \vspace{0.45cm}
    {\color{obeTeal}\rule{\textwidth}{1pt}}
    \vfill
    {\color{obeGray}\sffamily\small Dosen Pengampu\par}
    \vspace{0.08cm}
    {\sffamily\bfseries \DosenPengampu\par}
    \vspace{0.2cm}
    {\color{obeGray}\sffamily\small \SemesterTeks\par}
    \vspace{0.25cm}
\end{frame}
}

% Frame pemisah bagian
\newcommand{\PemisahBagian}[1]{%
\begin{frame}[plain]
    \vfill
    \begin{center}
        {\color{obeTeal}\sffamily\small\bfseries BAGIAN\par}
        \vspace{0.35cm}
        {\color{obeBlue}\sffamily\Large\bfseries #1\par}
        \vspace{0.35cm}
        {\color{obeTeal}\rule{0.55\textwidth}{0.9pt}}
    \end{center}
    \vfill
\end{frame}
}

% ---------------------------------------------------------
% 7. Gaya penampil kode
% ---------------------------------------------------------
% =========================================================
% ppt/src/listing-ppt.tex
% Gaya penampil kode untuk slide.
%
% Diturunkan dari buku_ajar/config/listings.tex, tetapi
% disesuaikan untuk slide: ukuran lebih kecil (karena frame
% 16:9 dan kode harus utuh terlihat) dan jarak lebih rapat.
% =========================================================

% ---------------------------------------------------------
% 1. Definisi bahasa MASM/TASM 8086 (salinan dari buku)
% ---------------------------------------------------------
\lstdefinelanguage{MASM}{
    morekeywords={
        MOV, XCHG, LEA, LDS, LES, PUSH, POP, PUSHA, POPA, PUSHF, POPF,
        LAHF, SAHF, XLAT, IN, OUT, CBW, CWD,
        ADD, ADC, SUB, SBB, INC, DEC, NEG, CMP, MUL, IMUL, DIV, IDIV,
        AAA, AAS, AAM, AAD, DAA, DAS,
        AND, OR, XOR, NOT, TEST, SHL, SHR, SAL, SAR, ROL, ROR, RCL, RCR,
        JMP, JE, JZ, JNE, JNZ, JG, JGE, JL, JLE, JA, JAE, JB, JBE,
        JC, JNC, JS, JNS, JO, JNO, JP, JPE, JNP, JPO, JCXZ,
        LOOP, LOOPE, LOOPZ, LOOPNE, LOOPNZ, CALL, RET, INT, IRET, INTO,
        HLT, NOP, CLI, STI, CLD, STD, CLC, STC, CMC, WAIT, ESC, LOCK,
        DB, DW, DD, DQ, DT, EQU, ORG, SEGMENT, ENDS, ASSUME, END,
        PROC, ENDP, MACRO, ENDM, LOCAL, PUBLIC, EXTRN, EXTERN,
        INCLUDE, INCLUDELIB, MODEL, STACK, CODE, DATA, CONST, ALIGN, EVEN,
        OFFSET, PTR, BYTE, WORD, DWORD, QWORD, TBYTE,
        NEAR, FAR, SHORT, LABEL, TYPE, SIZE, LENGTH, DUP, THIS,
        AX, BX, CX, DX, AL, AH, BL, BH, CL, CH, DL, DH,
        SP, BP, SI, DI, IP, CS, DS, SS, ES, FLAGS,
        CF, PF, AF, ZF, TF, IF, DF, OF
    },
    sensitive=false,
    morecomment=[l]{;},
    morestring=[b]',
    morestring=[b]"
}

% ---------------------------------------------------------
% 2. Gaya kode assembly untuk slide
% ---------------------------------------------------------
\lstdefinestyle{asmslide}{
    language=MASM,
    basicstyle=\ttfamily\scriptsize,
    keywordstyle=\color{obeBlue}\bfseries,
    commentstyle=\color{obeGreen}\itshape,
    stringstyle=\color{obeMaroon},
    numberstyle=\tiny\color{obeGray},
    numbers=left,
    numbersep=5pt,
    stepnumber=1,
    backgroundcolor=\color{codeBack},
    frame=single,
    rulecolor=\color{codeFrame},
    framesep=4pt,
    breaklines=true,
    breakatwhitespace=false,
    % CATATAN PENTING: opsi `postbreak` TIDAK boleh dipakai di sini.
    % Di beamer, `breaklines=true` bersama `postbreak=<kode>` membuat
    % listings gagal dirakit dengan galat
    %   "Missing number, treated as zero <to be read again> \kern".
    % Karena \lstset{style=asmslide} bersifat global, opsi postbreak
    % juga bocor ke gaya lain (mis. outslide/keluaran) dan merusaknya.
    % Di buku ajar (kelas `book`) justru sebaliknya: postbreak aman dan
    % dipakai. Jadi larangan ini khusus untuk slide.
    showstringspaces=false,
    tabsize=4,
    columns=fullflexible,
    keepspaces=true,
    upquote=true,
    xleftmargin=12pt,
    aboveskip=0.4em,
    belowskip=0.4em
}

% ---------------------------------------------------------
% 3. Gaya keluaran program / sesi debugger untuk slide
% ---------------------------------------------------------
\lstdefinestyle{outslide}{
    basicstyle=\ttfamily\scriptsize,
    keywordstyle=\color{black},
    commentstyle=\color{obeGray},
    stringstyle=\color{black},
    numbers=none,
    backgroundcolor=\color{obeGrayLight},
    frame=single,
    rulecolor=\color{codeFrame},
    framesep=4pt,
    breaklines=true,
    showstringspaces=false,
    columns=fullflexible,
    keepspaces=true,
    upquote=true,
    xleftmargin=4pt,
    aboveskip=0.4em,
    belowskip=0.4em
}

\lstset{style=asmslide}

% ---------------------------------------------------------
% 4. Environment ringkas untuk dipakai di dalam frame
%
% Frame yang memuat kedua environment ini WAJIB memakai
% opsi [fragile], karena isinya verbatim.
% ---------------------------------------------------------
\lstnewenvironment{kodeasm}[1][]{%
    \lstset{style=asmslide, #1}%
}{}

\lstnewenvironment{keluaran}[1][]{%
    \lstset{style=outslide, #1}%
}{}


%
% Palet warna disalin dari buku_ajar/config/styles.tex agar
% identitas visual slide sama dengan buku ajar.
% =========================================================

% ---------------------------------------------------------
% 1. Paket
% ---------------------------------------------------------
\usepackage[T1]{fontenc}
\usepackage[utf8]{inputenc}
\usepackage[indonesian]{babel}
\usepackage{lmodern}
\usepackage{tikz}
\usetikzlibrary{arrows.meta, positioning, calc, shapes.geometric}
\usepackage{booktabs}
\usepackage{array}
\usepackage{amsmath}
\usepackage{amsfonts}
\usepackage{pifont}          % \ding untuk tanda checklist
\usepackage{listings}

% ---------------------------------------------------------
% 2. Palet warna OBE (identik dengan buku)
% ---------------------------------------------------------
\definecolor{obeBlue}{HTML}{1F3864}
\definecolor{obeBlueLight}{HTML}{DCE4F5}
\definecolor{obeTeal}{HTML}{0F6E6E}
\definecolor{obeTealLight}{HTML}{D9F0EE}
\definecolor{obeGreen}{HTML}{2E6B34}
\definecolor{obeGreenLight}{HTML}{DEEFDF}
\definecolor{obeAmber}{HTML}{8A5A00}
\definecolor{obeAmberLight}{HTML}{FDF1D6}
\definecolor{obeMaroon}{HTML}{7A2233}
\definecolor{obeMaroonLight}{HTML}{F7E1E5}
\definecolor{obeGray}{HTML}{4A4A4A}
\definecolor{obeGrayLight}{HTML}{F0F0F0}
\definecolor{codeBack}{HTML}{F7F7F4}
\definecolor{codeFrame}{HTML}{C9C9C0}

% ---------------------------------------------------------
% 3. Identitas mata kuliah (dipakai seluruh deck)
% ---------------------------------------------------------
\newcommand{\NamaMK}{Bahasa Pemrograman Tingkat Rendah}
\newcommand{\NamaMKShort}{Bhs.\ Pemrograman Tingkat Rendah}
\newcommand{\KodeMK}{MWCA010}
\newcommand{\Prodi}{Program Studi Teknik Informatika}
\newcommand{\Fakultas}{Fakultas Teknologi Informasi}
\newcommand{\PerguruanTinggi}{Universitas Bale Bandung}
\newcommand{\DosenPengampu}{Cecep Suwanda, S.Si., M.Kom.}
\newcommand{\SemesterTeks}{Semester Ganjil}

% \PertemuanTeks dan \JudulBab ditetapkan di masing-masing
% berkas minggu-NN.tex sebelum \begin{document}.
\providecommand{\PertemuanTeks}{Pertemuan}

% ---------------------------------------------------------
% 4. Tema beamer
% ---------------------------------------------------------
\usetheme{default}
\setbeamertemplate{navigation symbols}{}

\setbeamercolor{normal text}{fg=black!88}
\setbeamercolor{structure}{fg=obeBlue}
\setbeamercolor{alerted text}{fg=obeMaroon}

% --- judul slide: pita biru penuh lebar ---
\setbeamercolor{frametitle}{fg=white,bg=obeBlue}
\setbeamerfont{frametitle}{series=\bfseries,size=\large}
\setbeamertemplate{frametitle}{%
    \nointerlineskip%
    \begin{beamercolorbox}[%
        wd=\paperwidth,ht=2.9ex,dp=1.5ex,
        leftskip=0.4cm,rightskip=0.4cm]{frametitle}%
        \usebeamerfont{frametitle}\insertframetitle%
    \end{beamercolorbox}%
    \nointerlineskip%
    \hbox{\textcolor{obeTeal}{\rule{\paperwidth}{1.1pt}}}%
}

% --- kaki slide: identitas mata kuliah, pertemuan, nomor ---
\setbeamercolor{footline}{fg=white,bg=obeGray}
\setbeamerfont{footline}{size=\scriptsize,family=\sffamily}
\setbeamertemplate{footline}{%
    \leavevmode%
    \hbox{%
        \begin{beamercolorbox}[wd=0.42\paperwidth,ht=2.1ex,dp=1.1ex,
            leftskip=0.4cm]{footline}%
            \usebeamerfont{footline}%
            \KodeMK\ \textbullet\ \NamaMKShort%
        \end{beamercolorbox}%
        \begin{beamercolorbox}[wd=0.32\paperwidth,ht=2.1ex,dp=1.1ex,
            center]{footline}%
            \usebeamerfont{footline}\PertemuanTeks%
        \end{beamercolorbox}%
        \begin{beamercolorbox}[wd=0.26\paperwidth,ht=2.1ex,dp=1.1ex,
            rightskip=0.4cm]{footline}%
            \usebeamerfont{footline}%
            \insertframenumber\,/\,\inserttotalframenumber%
        \end{beamercolorbox}%
    }%
    \vskip0pt%
}

% --- daftar ---
\setbeamercolor{itemize item}{fg=obeTeal}
\setbeamercolor{itemize subitem}{fg=obeTeal}
\setbeamercolor{enumerate item}{fg=obeBlue}
\setbeamertemplate{itemize item}{\raise0.2ex\hbox{\tiny$\blacksquare$}}
\setbeamertemplate{itemize subitem}{\raise0.2ex\hbox{\scriptsize$\bullet$}}

% --- kotak (pengganti tcolorbox pada buku) ---
\setbeamercolor{block title}{bg=obeBlueLight,fg=obeBlue}
\setbeamercolor{block body}{bg=obeBlueLight!45,fg=black!88}
\setbeamercolor{block title alerted}{bg=obeMaroonLight,fg=obeMaroon}
\setbeamercolor{block body alerted}{bg=obeMaroonLight!45,fg=black!88}
\setbeamercolor{block title example}{bg=obeTealLight,fg=obeTeal}
\setbeamercolor{block body example}{bg=obeTealLight!45,fg=black!88}

% ---------------------------------------------------------
% 5. Makro notasi (sama persis dengan buku)
%
% Catatan: \note TIDAK didefinisikan di sini karena beamer
% sudah memakai \note untuk catatan pembicara. Untuk catatan
% di slide dipakai \catatan.
% ---------------------------------------------------------
\newcommand{\code}[1]{\texttt{#1}}
\newcommand{\asm}[1]{\texttt{#1}}
\newcommand{\reg}[1]{\texttt{#1}}
\newcommand{\flag}[1]{\texttt{#1}}
\newcommand{\dir}[1]{\texttt{#1}}
\newcommand{\file}[1]{\texttt{#1}}
\newcommand{\eng}[1]{\textit{#1}}

\newcommand{\hex}[1]{\texttt{#1h}}
\newcommand{\bin}[1]{\texttt{#1b}}
\newcommand{\dec}[1]{\texttt{#1}}

\newcommand{\important}[1]{\textbf{Penting:} #1}
\newcommand{\catatan}[1]{\textbf{Catatan:} #1}
\newcommand{\term}[1]{\textbf{#1}}

\newcommand{\CPL}[1]{\textbf{CPL#1}}
\newcommand{\CPMK}[1]{\textbf{CPMK#1}}
\newcommand{\SubCPMK}[1]{\textbf{Sub-CPMK #1}}

\newcommand{\checkbox}{\ding{113}\hspace{0.35em}}
\newcommand{\checked}{\ding{51}\hspace{0.35em}}
\newcommand{\crossed}{\ding{55}\hspace{0.35em}}

\newcommand{\flowarrow}{\ensuremath{\rightarrow}}

% ---------------------------------------------------------
% 6. Makro tata letak slide
% ---------------------------------------------------------

% Keterangan sumber materi (ketertelusuran ke buku)
\newcommand{\Sumber}[1]{%
    \vfill
    {\scriptsize\color{obeGray}Sumber: #1\par}%
}

% Frame judul deck: \JudulDeck{<judul bab>}{<pertemuan>}
\newcommand{\JudulDeck}[2]{%
\begin{frame}[plain]
    \vspace{0.25cm}
    {\color{obeBlue}\sffamily\bfseries\small \NamaMK\ (\KodeMK)\par}
    \vspace{0.12cm}
    {\color{obeGray}\sffamily\scriptsize
     \Prodi\ \textbullet\ \Fakultas\ \textbullet\ \PerguruanTinggi\par}
    \vfill
    {\color{obeBlue}\sffamily\LARGE\bfseries #1\par}
    \vspace{0.4cm}
    {\color{obeTeal}\sffamily\large #2\par}
    \vspace{0.45cm}
    {\color{obeTeal}\rule{\textwidth}{1pt}}
    \vfill
    {\color{obeGray}\sffamily\small Dosen Pengampu\par}
    \vspace{0.08cm}
    {\sffamily\bfseries \DosenPengampu\par}
    \vspace{0.2cm}
    {\color{obeGray}\sffamily\small \SemesterTeks\par}
    \vspace{0.25cm}
\end{frame}
}

% Frame pemisah bagian
\newcommand{\PemisahBagian}[1]{%
\begin{frame}[plain]
    \vfill
    \begin{center}
        {\color{obeTeal}\sffamily\small\bfseries BAGIAN\par}
        \vspace{0.35cm}
        {\color{obeBlue}\sffamily\Large\bfseries #1\par}
        \vspace{0.35cm}
        {\color{obeTeal}\rule{0.55\textwidth}{0.9pt}}
    \end{center}
    \vfill
\end{frame}
}

% ---------------------------------------------------------
% 7. Gaya penampil kode
% ---------------------------------------------------------
% =========================================================
% ppt/src/listing-ppt.tex
% Gaya penampil kode untuk slide.
%
% Diturunkan dari buku_ajar/config/listings.tex, tetapi
% disesuaikan untuk slide: ukuran lebih kecil (karena frame
% 16:9 dan kode harus utuh terlihat) dan jarak lebih rapat.
% =========================================================

% ---------------------------------------------------------
% 1. Definisi bahasa MASM/TASM 8086 (salinan dari buku)
% ---------------------------------------------------------
\lstdefinelanguage{MASM}{
    morekeywords={
        MOV, XCHG, LEA, LDS, LES, PUSH, POP, PUSHA, POPA, PUSHF, POPF,
        LAHF, SAHF, XLAT, IN, OUT, CBW, CWD,
        ADD, ADC, SUB, SBB, INC, DEC, NEG, CMP, MUL, IMUL, DIV, IDIV,
        AAA, AAS, AAM, AAD, DAA, DAS,
        AND, OR, XOR, NOT, TEST, SHL, SHR, SAL, SAR, ROL, ROR, RCL, RCR,
        JMP, JE, JZ, JNE, JNZ, JG, JGE, JL, JLE, JA, JAE, JB, JBE,
        JC, JNC, JS, JNS, JO, JNO, JP, JPE, JNP, JPO, JCXZ,
        LOOP, LOOPE, LOOPZ, LOOPNE, LOOPNZ, CALL, RET, INT, IRET, INTO,
        HLT, NOP, CLI, STI, CLD, STD, CLC, STC, CMC, WAIT, ESC, LOCK,
        DB, DW, DD, DQ, DT, EQU, ORG, SEGMENT, ENDS, ASSUME, END,
        PROC, ENDP, MACRO, ENDM, LOCAL, PUBLIC, EXTRN, EXTERN,
        INCLUDE, INCLUDELIB, MODEL, STACK, CODE, DATA, CONST, ALIGN, EVEN,
        OFFSET, PTR, BYTE, WORD, DWORD, QWORD, TBYTE,
        NEAR, FAR, SHORT, LABEL, TYPE, SIZE, LENGTH, DUP, THIS,
        AX, BX, CX, DX, AL, AH, BL, BH, CL, CH, DL, DH,
        SP, BP, SI, DI, IP, CS, DS, SS, ES, FLAGS,
        CF, PF, AF, ZF, TF, IF, DF, OF
    },
    sensitive=false,
    morecomment=[l]{;},
    morestring=[b]',
    morestring=[b]"
}

% ---------------------------------------------------------
% 2. Gaya kode assembly untuk slide
% ---------------------------------------------------------
\lstdefinestyle{asmslide}{
    language=MASM,
    basicstyle=\ttfamily\scriptsize,
    keywordstyle=\color{obeBlue}\bfseries,
    commentstyle=\color{obeGreen}\itshape,
    stringstyle=\color{obeMaroon},
    numberstyle=\tiny\color{obeGray},
    numbers=left,
    numbersep=5pt,
    stepnumber=1,
    backgroundcolor=\color{codeBack},
    frame=single,
    rulecolor=\color{codeFrame},
    framesep=4pt,
    breaklines=true,
    breakatwhitespace=false,
    % CATATAN PENTING: opsi `postbreak` TIDAK boleh dipakai di sini.
    % Di beamer, `breaklines=true` bersama `postbreak=<kode>` membuat
    % listings gagal dirakit dengan galat
    %   "Missing number, treated as zero <to be read again> \kern".
    % Karena \lstset{style=asmslide} bersifat global, opsi postbreak
    % juga bocor ke gaya lain (mis. outslide/keluaran) dan merusaknya.
    % Di buku ajar (kelas `book`) justru sebaliknya: postbreak aman dan
    % dipakai. Jadi larangan ini khusus untuk slide.
    showstringspaces=false,
    tabsize=4,
    columns=fullflexible,
    keepspaces=true,
    upquote=true,
    xleftmargin=12pt,
    aboveskip=0.4em,
    belowskip=0.4em
}

% ---------------------------------------------------------
% 3. Gaya keluaran program / sesi debugger untuk slide
% ---------------------------------------------------------
\lstdefinestyle{outslide}{
    basicstyle=\ttfamily\scriptsize,
    keywordstyle=\color{black},
    commentstyle=\color{obeGray},
    stringstyle=\color{black},
    numbers=none,
    backgroundcolor=\color{obeGrayLight},
    frame=single,
    rulecolor=\color{codeFrame},
    framesep=4pt,
    breaklines=true,
    showstringspaces=false,
    columns=fullflexible,
    keepspaces=true,
    upquote=true,
    xleftmargin=4pt,
    aboveskip=0.4em,
    belowskip=0.4em
}

\lstset{style=asmslide}

% ---------------------------------------------------------
% 4. Environment ringkas untuk dipakai di dalam frame
%
% Frame yang memuat kedua environment ini WAJIB memakai
% opsi [fragile], karena isinya verbatim.
% ---------------------------------------------------------
\lstnewenvironment{kodeasm}[1][]{%
    \lstset{style=asmslide, #1}%
}{}

\lstnewenvironment{keluaran}[1][]{%
    \lstset{style=outslide, #1}%
}{}



\renewcommand{\PertemuanTeks}{Pertemuan 12}

\begin{document}

% =========================================================
% 1. JUDUL
% =========================================================
\JudulDeck{Prosedur dan Manajemen Stack}{Pertemuan 12}

% =========================================================
% 2. CAPAIAN DAN TUJUAN
% =========================================================
\begin{frame}{Capaian dan Tujuan Pembelajaran}
\begin{block}{Sub-CPMK 3.4}
\small Mengembangkan prosedur/subrutin dengan mekanisme pengiriman
parameter (\eng{parameter passing}) dan mengelola \eng{stack frame}
melalui instruksi \asm{PUSH} dan \asm{POP}, register \reg{BP}, serta
konvensi pemanggilan yang benar.
\end{block}

{\small\textbf{CPMK yang didukung:} \CPMK{03}\quad\textbullet\quad
\textbf{CPL yang didukung:} \CPL{03}\par}

\begin{exampleblock}{Setelah pertemuan ini, mahasiswa mampu:}
\footnotesize
\begin{enumerate}
    \item menjelaskan alasan pemakaian prosedur beserta biaya
          pemanggilannya;
    \item menjelaskan kerja \asm{CALL} dan \asm{RET} serta membedakan
          pemanggilan \eng{near} dan \eng{far};
    \item menelusuri isi stack sel per sel pada \asm{PUSH}, \asm{POP},
          dan \asm{CALL};
    \item membangun \eng{stack frame} melalui \eng{prologue} dan
          \eng{epilogue} yang benar;
    \item mengirim parameter melalui register, stack, dan variabel
          global beserta kelebihan serta kekurangannya;
    \item menulis prosedur rekursif dan menjelaskan risiko
          \eng{stack overflow}.
\end{enumerate}
\end{exampleblock}
\end{frame}

% =========================================================
% 3. PETA MATERI
% =========================================================
\begin{frame}{Peta Materi}
\small
\begin{enumerate}
    \item \textbf{Prosedur, Pemanggilan, dan Pengembalian}\\
          {\footnotesize alasan prosedur; biaya pemanggilan;
          \asm{CALL} dan \asm{RET}; alamat kembali \reg{CS}:\reg{IP};
          pemanggilan \eng{near} dan \eng{far}}

    \vspace{0.25em}
    \item \textbf{Stack dan Operasi Stack}\\
          {\footnotesize LIFO; \reg{SS}:\reg{SP}; arah tumbuh ke
          alamat rendah; \asm{PUSH}, \asm{POP}, \asm{PUSHF},
          \asm{POPF}; penelusuran isi stack; \reg{BP} penunjuk
          bingkai}

    \vspace{0.25em}
    \item \textbf{Stack Frame, Konvensi Pemanggilan, dan Rekursi}\\
          {\footnotesize \eng{prologue} dan \eng{epilogue}; akses
          parameter; konvensi pemanggilan; nilai balik; rekursi dan
          \eng{stack overflow}}
\end{enumerate}

\vspace{0.35em}
\begin{block}{Praktik pada pertemuan ini}
\small Menulis prosedur berparameter dan memeriksa isi stack dengan
\eng{debugger} pada \file{prosedur-luas.asm} dan \file{faktorial.asm}.
\end{block}
\end{frame}

% =========================================================
% BAGIAN 1
% =========================================================
\PemisahBagian{Prosedur, Pemanggilan, dan Pengembalian}

% ---------------------------------------------------------
\begin{frame}{Mengapa Prosedur Diperlukan}
\small
Program yang ditulis sebagai satu deretan instruksi panjang sulit
dibaca, diuji, dan diperbaiki. Solusinya adalah \term{prosedur}.

\vspace{0.3em}
\begin{enumerate}
    \item \term{Modularitas.} Setiap prosedur menangani satu tugas
          yang batasnya jelas, sehingga program besar menjadi
          kumpulan bagian kecil yang dapat diuji terpisah.

    \vspace{0.2em}
    \item \term{Penggunaan ulang} (\eng{reuse}). Kode yang sering
          dipakai cukup ditulis sekali, misalnya prosedur pencetak
          bilangan heksadesimal.

    \vspace{0.2em}
    \item \term{Abstraksi.} Pemanggil cukup mengetahui antarmuka
          prosedur: nama, parameter, dan nilai baliknya. Rincian di
          dalamnya dapat diganti tanpa mengubah pemanggil.
\end{enumerate}

\vspace{0.25em}
\begin{block}{Prosedur dan \asm{JMP}}
\small \asm{JMP} hanya memindahkan kendali tanpa menyediakan cara
kembali; prosedur melengkapinya dengan mekanisme pengembalian.
\end{block}
\Sumber{Bab XIII, Mengapa Prosedur Diperlukan}
\end{frame}

% ---------------------------------------------------------
\begin{frame}{Definisi dan Biaya Pemanggilan Prosedur}
\small
\begin{exampleblock}{Definisi}
\small \term{Prosedur} (disebut juga \term{subrutin}) adalah
sekelompok instruksi bernama yang dapat dipanggil dari bagian lain
program, menerima parameter, dan dapat mengembalikan nilai kepada
pemanggilnya.
\end{exampleblock}

\vspace{0.15em}
{\small Pemakaian prosedur tidak gratis. Setiap pemanggilan
menimbulkan \term{biaya pemanggilan prosedur}:\par}

\vspace{0.12em}
{\footnotesize
\begin{center}
\begin{tabular}{@{}p{5.4cm}cp{6.0cm}@{}}
    \toprule
    \textbf{Tahap} & \textbf{Siklus} & \textbf{Keterangan} \\
    \midrule
    \asm{PUSH} parameter & 11--15 &
    setiap parameter menimbulkan akses memori \\
    \addlinespace
    \asm{CALL} prosedur & 19 &
    menyimpan alamat kembali lalu melompat \\
    \addlinespace
    \asm{RET} & 16 &
    mengambil alamat kembali dari stack \\
    \bottomrule
\end{tabular}
\end{center}}

\vspace{0.2em}
{\footnotesize\color{obeGray}Untuk prosedur yang sangat sederhana,
biaya ini dapat melebihi biaya pekerjaannya sendiri. Karena itu
prosedur dipakai bukan untuk mengejar kecepatan, melainkan untuk
kejelasan dan kemudahan pemeliharaan.\par}
\Sumber{Bab XIII, Mengapa Prosedur Diperlukan}
\end{frame}

% ---------------------------------------------------------
\begin{frame}[fragile]{Instruksi \asm{CALL} dan \asm{RET}}
\small
\asm{CALL} mengerjakan dua hal sekaligus: menyimpan \term{alamat
kembali} --- nilai \reg{CS} dan \reg{IP} pada instruksi \emph{sesudah}
\asm{CALL} --- lalu memuat alamat prosedur ke \reg{CS} dan \reg{IP}.

\vspace{0.25em}
\begin{kodeasm}[numbers=none]
    CALL HitungLuas        ; simpan alamat kembali, lalu lompat
    MOV  luas, AX          ; dieksekusi setelah RET
\end{kodeasm}

\vspace{0.25em}
{\small
\begin{itemize}
    \item Alamat kembali yang disimpan adalah alamat \code{MOV luas,
          AX}; \asm{RET} memuatnya kembali ke \reg{IP}.
    \item Stack dengan demikian merupakan bagian dari mekanisme
          kendali program, bukan sekadar tempat menyimpan data.
\end{itemize}}

\vspace{0.2em}
\begin{alertblock}{\asm{CALL} tidak sama dengan \asm{JMP}}
\small \asm{JMP} tidak menyimpan alamat kembali, sehingga \asm{RET}
akan mengambil alamat yang salah dari stack.
\end{alertblock}
\end{frame}

% ---------------------------------------------------------
\begin{frame}{Pemanggilan \eng{Near} dan \eng{Far}}
\footnotesize
\begin{center}
\begin{tabular}{@{}p{3.0cm}p{4.9cm}p{4.9cm}@{}}
    \toprule
    \textbf{Aspek} & \textbf{\asm{CALL} \eng{near}} &
    \textbf{\asm{CALL} \eng{far}} \\
    \midrule
    Isi yang didorong ke stack & \reg{IP} saja (\dec{2} byte) &
    \reg{CS} lalu \reg{IP} (\dec{4} byte) \\
    \addlinespace
    Instruksi kembali & \asm{RET} (\code{C3h}) &
    \asm{RETF} (\code{CBh}) \\
    \addlinespace
    Bentuk penulisan & \asm{CALL HitungLuas} &
    \asm{CALL LUAR:HitungLuas} \\
    \addlinespace
    Pemakaian lazim & Prosedur pada satu berkas program &
    Prosedur pada pustaka atau \eng{overlay} lain \\
    \bottomrule
\end{tabular}
\end{center}

\vspace{0.35em}
{\footnotesize Pada \term{pemanggilan dekat} (\eng{near call})
prosedur berada pada segmen kode yang sama, sehingga hanya \reg{IP}
yang perlu disimpan. Pada \term{pemanggilan jauh} (\eng{far call})
prosedur berada pada segmen lain, sehingga \reg{CS} dan \reg{IP}
sama-sama disimpan. Seluruh contoh pada bab ini memakai pemanggilan
\eng{near}.\par}
\Sumber{Bab XIII, Pemanggilan Near dan Far}
\end{frame}

% ---------------------------------------------------------
\begin{frame}{Rangkuman Bagian 1}
\small
\begin{itemize}
    \item Prosedur dipakai untuk memperoleh modularitas, penggunaan
          ulang, dan abstraksi; \asm{JMP} tidak dapat menggantikannya
          karena tidak menyediakan cara kembali.

    \vspace{0.25em}
    \item Biaya pemanggilan mencakup dorongan parameter
          (\dec{11}--\dec{15} siklus per \asm{PUSH}), \asm{CALL}
          (\dec{19} siklus), dan \asm{RET} (\dec{16} siklus) ---
          sehingga prosedur dipakai demi kejelasan, bukan kecepatan.

    \vspace{0.25em}
    \item \asm{CALL} menyimpan alamat kembali (\reg{CS}:\reg{IP} pada
          instruksi sesudah \asm{CALL}) ke stack, lalu melompat;
          \asm{RET} mengambil alamat itu kembali ke \reg{IP}.

    \vspace{0.25em}
    \item Pemanggilan \eng{near} menyimpan \reg{IP} saja
          (\dec{2} byte) dan ditutup \asm{RET}; pemanggilan \eng{far}
          menyimpan \reg{CS} dan \reg{IP} (\dec{4} byte) dan ditutup
          \asm{RETF}.
\end{itemize}
\end{frame}

% =========================================================
% BAGIAN 2
% =========================================================
\PemisahBagian{Stack dan Operasi Stack}

% ---------------------------------------------------------
\begin{frame}{Stack: LIFO dan Arah Tumbuh}
\small
\begin{itemize}
    \item \term{Stack} adalah daerah memori yang dipakai dengan aturan
          \eng{last in, first out} (LIFO): data yang terakhir
          dimasukkan adalah data yang pertama dikeluarkan.

    \vspace{0.2em}
    \item Stack 8086 ditunjuk oleh pasangan \reg{SS} dan \reg{SP};
          alamat fisik sel teratas adalah
          \reg{SS} $\times$ \hex{10} $+$ \reg{SP}.

    \vspace{0.2em}
    \item \important{Setiap penambahan isi stack menurunkan nilai
          \reg{SP}, sehingga stack tumbuh menuju alamat yang semakin
          rendah.}
\end{itemize}

\vspace{0.3em}
\begin{exampleblock}{Contoh angka}
\small Bila \reg{SS} $=$ \hex{1000} dan ruang stack \hex{100}, maka
\reg{SP} mula-mula bernilai \hex{FFFE} --- alamat tertinggi dari ruang
yang disediakan direktif \dir{.STACK} --- dan alamat fisik sel
teratasnya adalah \hex{1FFFE}.
\end{exampleblock}
\Sumber{Bab XIII, Pengertian dan Arah Tumbuh Stack}
\end{frame}

% ---------------------------------------------------------
\begin{frame}[fragile]{Instruksi \asm{PUSH} dan \asm{POP}}
\small
\begin{kodeasm}[numbers=none]
PUSH sumber    ; SP <- SP - 2, lalu isi sumber ke memori SS:SP
POP  tujuan    ; isi SS:SP ke tujuan, lalu SP <- SP + 2
\end{kodeasm}

\vspace{0.2em}
\small
\begin{itemize}
    \item Keduanya selalu bekerja pada satu \term{kata}
          (\eng{word}) berukuran \dec{2} byte.
    \item Pada 8086, \asm{PUSH} hanya menerima register atau isi
          memori 16-bit, sehingga \asm{PUSH 5} tidak sah.
    \item \asm{PUSHA} dan \asm{POPA} belum tersedia karena baru
          diperkenalkan pada prosesor 80186.
\end{itemize}

\vspace{0.2em}
\begin{exampleblock}{\asm{PUSHF} dan \asm{POPF}}
\small \asm{PUSHF} mendorong seluruh isi register \reg{FLAGS} sebagai
satu kata 16-bit, \asm{POPF} memulihkannya. Dipakai bila prosedur
mengubah \eng{flag} untuk keperluan internal. Perlu diingat bahwa
\asm{POPF} ikut memulihkan \flag{IF}, sehingga pemulihan yang salah
dapat mematikan interupsi.
\end{exampleblock}
\end{frame}

% ---------------------------------------------------------
\begin{frame}{Jejak \asm{PUSH AX} pada Stack}
\begin{center}
\begin{tikzpicture}[
    sel/.style={draw=obeBlue,line width=0.6pt,minimum width=1.9cm,
        minimum height=0.5cm,font=\ttfamily\footnotesize,inner sep=1pt},
    alamat/.style={font=\ttfamily\scriptsize,text=obeGray,anchor=west},
    judul/.style={font=\scriptsize\bfseries,text=obeBlue},
    penunjuk/.style={-{Stealth[length=2.2mm]},line width=1pt,obeTeal}
]
    % --- keadaan sebelum ---
    \node[judul] (ja) at (0, 1.10) {sebelum \asm{PUSH AX}};
    \node[sel,fill=obeGrayLight] (a0) at (0, 0) {\textit{bebas}};
    \node[alamat] at ([xshift=0.12cm]a0.east) {\hex{FFFE}};
    \node[font=\ttfamily\scriptsize,text=obeBlue,left=0.5cm of a0] (spa)
        {SP};
    \draw[penunjuk] (spa.east) -- (a0.west);

    % --- keadaan sesudah ---
    \node[judul] (jb) at (5.0, 1.10) {setelah \asm{PUSH AX}};
    \node[sel,fill=obeGrayLight] (b0) at (5.0, 0.55) {\textit{bebas}};
    \node[sel,fill=obeAmberLight] (b1) at (5.0, 0) {1234h};
    \node[alamat] at ([xshift=0.12cm]b0.east) {\hex{FFFE}};
    \node[alamat] at ([xshift=0.12cm]b1.east) {\hex{FFFC}};
    \node[font=\ttfamily\scriptsize,text=obeBlue,left=0.5cm of b1] (spb)
        {SP};
    \draw[penunjuk] (spb.east) -- (b1.west);
\end{tikzpicture}
\end{center}

\vspace{0.3em}
{\footnotesize\color{obeGray}Keadaan awal \reg{SP} $=$ \hex{FFFE} dan
\reg{AX} $=$ \hex{1234}. Setelah \asm{PUSH AX}, \reg{SP} turun ke
\hex{FFFC}; byte rendah \hex{34} ditulis pada \hex{FFFC} dan byte
tinggi \hex{12} pada \hex{FFFD} --- byte rendah selalu berada pada
alamat yang lebih rendah.\par}
\Sumber{Bab XIII, Instruksi PUSH dan POP}
\end{frame}

% ---------------------------------------------------------
\begin{frame}{Menelusuri Isi Stack pada Pemanggilan (1/2)}
\footnotesize Keadaan awal: \reg{SS} $=$ \hex{1000}, \reg{SP} $=$
\hex{FFFE}, stack kosong.

\vspace{0.15em}
\begin{center}
\begin{tabular}{@{}p{0.6cm}p{2.8cm}p{1.65cm}p{2.35cm}p{3.9cm}@{}}
    \toprule
    \textbf{\#} & \textbf{Instruksi} & \textbf{\reg{SP} sesudah} &
    \textbf{Isi sel teratas} & \textbf{Keterangan} \\
    \midrule
    1 & (keadaan awal) & \hex{FFFE} & kosong &
    \reg{SP} menunjuk sel teratas \\
    \addlinespace
    2 & \asm{PUSH} \code{tinggi} & \hex{FFFC} & \hex{000A} &
    parameter kedua masuk lebih dahulu \\
    \addlinespace
    3 & \asm{PUSH} \code{lebar} & \hex{FFFA} & \hex{0005} &
    parameter pertama masuk paling akhir \\
    \addlinespace
    4 & \asm{CALL} \code{HitungLuas} & \hex{FFF8} & \hex{0102} &
    alamat kembali (\reg{IP}) didorong \\
    \bottomrule
\end{tabular}
\end{center}

\vspace{0.25em}
{\footnotesize\color{obeGray}Urutan sel itu bukan kebetulan,
melainkan akibat langsung dari urutan \asm{PUSH} dan \asm{CALL}:
parameter yang didorong paling akhir berada paling dekat dengan
\reg{BP}.\par}
\Sumber{Bab XIII, Menelusuri Isi Stack pada Pemanggilan Prosedur}
\end{frame}

% ---------------------------------------------------------
\begin{frame}{Menelusuri Isi Stack pada Pemanggilan (2/2)}
\footnotesize Lanjutan jejak: \eng{prologue}, \eng{epilogue}, dan
\asm{RET}.

\vspace{0.15em}
\begin{center}
\begin{tabular}{@{}p{0.6cm}p{2.8cm}p{1.65cm}p{2.35cm}p{3.9cm}@{}}
    \toprule
    \textbf{\#} & \textbf{Instruksi} & \textbf{\reg{SP} sesudah} &
    \textbf{Isi sel teratas} & \textbf{Keterangan} \\
    \midrule
    5 & \asm{PUSH} \reg{BP} & \hex{FFF6} & \hex{0000} &
    \reg{BP} milik pemanggil disimpan \\
    \addlinespace
    6 & \asm{MOV} \reg{BP}, \reg{SP} & \hex{FFF6} & \hex{0000} &
    \reg{BP} $=$ \hex{FFF6}; bingkai stack siap \\
    \addlinespace
    7 & \asm{POP} \reg{BP} & \hex{FFF8} & \hex{0102} &
    \reg{BP} milik pemanggil dipulihkan \\
    \addlinespace
    8 & \asm{RET} \code{4} & \hex{FFFE} & kosong &
    \reg{IP} diambil, lalu \reg{SP} naik \dec{4} byte \\
    \bottomrule
\end{tabular}
\end{center}

\vspace{0.25em}
{\footnotesize\color{obeGray}Langkah 1--4 terlihat pada frame
sebelumnya. Perhatikan bahwa \asm{MOV BP, SP} tidak mengubah
\reg{SP}: yang berubah hanyalah nilai \reg{BP} sebagai titik acuan
bingkai.\par}
\Sumber{Bab XIII, Menelusuri Isi Stack pada Pemanggilan Prosedur}
\end{frame}

% ---------------------------------------------------------
\begin{frame}{Isi Stack Sel per Sel Setelah \eng{Prologue}}
\begin{center}
\begin{tikzpicture}[
    sel/.style={draw=obeBlue,line width=0.6pt,minimum width=2.9cm,
        minimum height=0.55cm,font=\ttfamily\small,inner sep=1pt},
    alamat/.style={font=\ttfamily\scriptsize,text=obeGray,anchor=west},
    catatan/.style={font=\scriptsize,text=obeMaroon,anchor=west},
    penunjuk/.style={-{Stealth[length=2.4mm]},line width=1pt,obeBlue}
]
    \node[sel,fill=obeGrayLight] (c0) {\textit{(bebas)}};
    \node[sel,fill=obeAmberLight,below=0pt of c0] (c1) {000Ah};
    \node[sel,fill=obeAmberLight,below=0pt of c1] (c2) {0005h};
    \node[sel,fill=obeBlueLight,below=0pt of c2] (c3) {0102h};
    \node[sel,fill=obeGreenLight,below=0pt of c3] (c4) {0000h};
    \node[sel,fill=obeGrayLight,below=0pt of c4] (c5) {\textit{(bebas)}};

    \node[alamat] at ([xshift=0.12cm]c0.east) {\hex{FFFE}};
    \node[alamat] at ([xshift=0.12cm]c1.east) {\hex{FFFC}};
    \node[alamat] at ([xshift=0.12cm]c2.east) {\hex{FFFA}};
    \node[alamat] at ([xshift=0.12cm]c3.east) {\hex{FFF8}};
    \node[alamat] at ([xshift=0.12cm]c4.east) {\hex{FFF6}};
    \node[alamat] at ([xshift=0.12cm]c5.east) {\hex{FFF4}};

    \node[catatan] at ([xshift=1.25cm]c1.east)
        {parameter 2: tinggi \code{[BP+6]}};
    \node[catatan] at ([xshift=1.25cm]c2.east)
        {parameter 1: lebar \code{[BP+4]}};
    \node[catatan] at ([xshift=1.25cm]c3.east)
        {alamat kembali \code{[BP+2]}};
    \node[catatan] at ([xshift=1.25cm]c4.east)
        {\reg{BP} lama \code{[BP+0]}};
    \node[catatan] at ([xshift=1.25cm]c5.east)
        {di bawah \reg{SP}: ruang bebas};

    \node[font=\scriptsize,text=obeTeal,above=0.06cm of c0]
        {alamat tinggi};
    \node[font=\scriptsize,text=obeTeal,below=0.06cm of c5]
        {alamat rendah};

    \node[font=\ttfamily\scriptsize,text=obeBlue,left=0.75cm of c4] (sp)
        {SP = BP};
    \draw[penunjuk] (sp.east) -- (c4.west);
\end{tikzpicture}
\end{center}

\vspace{0.2em}
{\footnotesize\color{obeGray}Keadaan pada langkah keenam. Setiap sel
berukuran \dec{2} byte.\par}
\Sumber{Bab XIII, Menelusuri Isi Stack pada Pemanggilan Prosedur}
\end{frame}

% ---------------------------------------------------------
\begin{frame}{Register \reg{BP} sebagai Penunjuk Bingkai}
\small
\begin{itemize}
    \item Akses parameter tidak dilakukan langsung melalui \reg{SP}
          karena \reg{SP} terus berubah setiap kali prosedur
          mendorong atau mengambil data.

    \vspace{0.2em}
    \item Register \reg{BP} (\eng{base pointer}) menyimpan salinan
          \reg{SP} melalui \asm{MOV BP, SP} dan tidak diubah oleh
          operasi \asm{PUSH} maupun \asm{POP} berikutnya.

    \vspace{0.2em}
    \item Karena itu jarak \reg{BP} ke setiap parameter tetap
          konstan. Pola inilah yang disebut \term{bingkai stack}
          (\eng{stack frame}).
\end{itemize}

\vspace{0.3em}
\begin{block}{Segmen bawaan \reg{BP} adalah \reg{SS}}
\small \code{[BP+4]} selalu mengacu ke stack tanpa perlu awalan
segmen. Cara ini merupakan penerapan mode pengalamatan
\eng{base-relative} yang dibahas pada Bab VIII.
\end{block}
\Sumber{Bab XIII, Register BP sebagai Penunjuk Bingkai}
\end{frame}

% ---------------------------------------------------------
\begin{frame}{Rangkuman Bagian 2}
\small
\begin{itemize}
    \item Stack bekerja secara LIFO dan ditunjuk \reg{SS}:\reg{SP};
          alamat fisik sel teratas $=$ \reg{SS} $\times$ \hex{10}
          $+$ \reg{SP}.

    \vspace{0.25em}
    \item Stack 8086 tumbuh ke arah alamat rendah: \asm{PUSH}
          menurunkan \reg{SP} \dec{2} byte lebih dahulu lalu menyimpan
          data, sedangkan \asm{POP} mengambil data lebih dahulu lalu
          menaikkan \reg{SP}.

    \vspace{0.25em}
    \item Setiap operasi bekerja pada satu kata (\dec{2} byte); byte
          rendah selalu berada pada alamat yang lebih rendah.
          \asm{PUSHF} dan \asm{POPF} menyimpan serta memulihkan
          register \reg{FLAGS}.

    \vspace{0.25em}
    \item Penelusuran satu pemanggilan lengkap: dua \asm{PUSH}
          parameter, \asm{CALL}, \asm{PUSH BP}, \asm{MOV BP, SP},
          lalu \asm{POP BP} dan \asm{RET 4} mengembalikan \reg{SP} ke
          \hex{FFFE}.

    \vspace{0.25em}
    \item \reg{BP} dipakai sebagai titik acuan karena nilainya tidak
          berubah oleh \asm{PUSH}/\asm{POP}, sehingga jarak ke setiap
          parameter tetap.
\end{itemize}
\end{frame}

% =========================================================
% BAGIAN 3
% =========================================================
\PemisahBagian{Stack Frame, Konvensi Pemanggilan, dan Rekursi}

% ---------------------------------------------------------
\begin{frame}[fragile]{\eng{Prologue} dan \eng{Epilogue}}
\small Prosedur selalu mengerjakan empat hal: membangun bingkai,
mengakses parameter, menghasilkan nilai balik, lalu membongkar
bingkai. Pola bakunya disebut \term{stack frame} standar.

\vspace{0.2em}
\begin{kodeasm}[numbers=none]
HitungLuas PROC
    PUSH BP            ; prologue : simpan BP pemanggil
    MOV  BP, SP        ; prologue : BP <- titik acuan bingkai

    ; --- badan prosedur ---
    ; parameter diakses sebagai [BP+4], [BP+6], ...

    POP  BP            ; epilogue : pulihkan BP pemanggil
    RET                ; epilogue : kembali ke pemanggil
HitungLuas ENDP
\end{kodeasm}

\vspace{0.15em}
{\footnotesize\color{obeGray}\asm{PUSH BP} wajib ada karena \reg{BP}
hanya satu buah; tanpa menyimpannya, prosedur akan menimpa \reg{BP}
milik pemanggil. Pasangan \asm{PUSH BP} dan \asm{POP BP} harus
seimbang pada setiap jalur keluar prosedur.\par}
\Sumber{Bab XIII, Prologue dan Epilogue}
\end{frame}

% ---------------------------------------------------------
\begin{frame}[fragile]{Mengakses Parameter dan Variabel Lokal}
\begin{columns}[t]
\begin{column}{0.5\textwidth}
\begin{kodeasm}[numbers=none]
    PUSH BP              ; prologue
    MOV  BP, SP
    SUB  SP, 4           ; dua var lokal
    MOV  WORD PTR [BP-2], 0
    MOV  WORD PTR [BP-4], 0
    MOV  SP, BP          ; epilogue
    POP  BP
    RET
\end{kodeasm}
\end{column}
\begin{column}{0.46\textwidth}
{\footnotesize\bfseries\color{obeBlue}Susunan bingkai}\\[0.2em]
{\footnotesize
\begin{itemize}
    \item \code{[BP+0]} --- \reg{BP} milik pemanggil
    \item \code{[BP+2]} --- alamat kembali (\reg{IP})
    \item \code{[BP+4]} --- parameter pertama
    \item \code{[BP+6]} --- parameter kedua
    \item \code{[BP+8]} --- parameter ketiga
    \item \code{[BP-2]} --- variabel lokal pertama
\end{itemize}}
\end{column}
\end{columns}

\vspace{0.2em}
\begin{alertblock}{Urutan \asm{PUSH} menentukan nomor offset}
\footnotesize Karena pemanggil mendorong parameter kedua lebih
dahulu, parameter pertama-lah yang berada pada \code{[BP+4]}. Bila
urutannya dibalik tanpa mengubah badan prosedur, kedua nilai akan
tertukar tanpa pesan kesalahan apa pun.
\end{alertblock}
\end{frame}

% ---------------------------------------------------------
\begin{frame}{Konvensi Pemanggilan: Tiga Cara Pengiriman Parameter}
\footnotesize
\begin{center}
\begin{tabular}{@{}p{2.1cm}p{3.3cm}p{3.5cm}p{3.2cm}@{}}
    \toprule
    \textbf{Cara} & \textbf{Kelebihan} & \textbf{Kekurangan} &
    \textbf{Contoh} \\
    \midrule
    Melalui register & tidak ada akses memori; paling cepat dan
    ringkas & jumlah parameter terbatas oleh jumlah register &
    \asm{MOV AX, 5} lalu \asm{CALL Tulis} \\
    \addlinespace
    Melalui stack & jumlah parameter tidak terbatas; pemanggil tidak
    perlu tahu register yang dipakai prosedur & lebih lambat karena
    setiap parameter menimbulkan akses memori; rawan stack tidak
    seimbang & \asm{PUSH 5} lalu \asm{CALL HitungLuas} \\
    \addlinespace
    Melalui variabel global & prosedur sangat sederhana; tidak ada
    instruksi tambahan & prosedur tidak dapat dipakai ulang; sulit
    ditelusuri saat kesalahan terjadi & \asm{MOV lebar, 5} lalu
    \asm{CALL HitungLuas} \\
    \bottomrule
\end{tabular}
\end{center}

\vspace{0.3em}
{\footnotesize\color{obeGray}\term{Konvensi pemanggilan} (\eng{calling
convention}) adalah kesepakatan tentang tiga hal: letak parameter,
letak nilai balik, dan siapa yang membersihkan stack. Tanpa
kesepakatan ini, prosedur yang ditulis satu orang tidak dapat dipakai
oleh program orang lain.\par}
\Sumber{Bab XIII, Konvensi Pemanggilan}
\end{frame}

% ---------------------------------------------------------
\begin{frame}[fragile]{Siapa yang Membersihkan Stack}
\small
\begin{columns}[t]
\begin{column}{0.46\textwidth}
\begin{exampleblock}{Prosedur membersihkan}
\footnotesize \eng{Epilogue} berakhir dengan \asm{RET} \code{n}:
mengambil \reg{IP} lalu menambahkan \code{n} pada \reg{SP} sekaligus.
\end{exampleblock}
\begin{kodeasm}[numbers=none]
    POP  BP
    RET  4        ; SP <- SP + 4
\end{kodeasm}
\end{column}
\begin{column}{0.48\textwidth}
\begin{exampleblock}{Pemanggil membersihkan}
\footnotesize Prosedur memakai \asm{RET} biasa, dan pemanggil
menambahkan \asm{ADD SP, 4} setelah \asm{CALL}.
\end{exampleblock}
\begin{kodeasm}[numbers=none]
    CALL HitungLuas
    ADD  SP, 4    ; buang argumen
\end{kodeasm}
\end{column}
\end{columns}

\vspace{0.25em}
\begin{alertblock}{Jangan dipakai bersama-sama}
\small \asm{RET} \code{4} yang disusul \asm{ADD SP, 4} membuat
\reg{SP} naik \dec{8} byte padahal hanya \dec{4} byte yang perlu
dibuang.
\end{alertblock}
\Sumber{Bab XIII, Konvensi Pemanggilan}
\end{frame}

% ---------------------------------------------------------
\begin{frame}{Tanggung Jawab Penyimpanan Register}
\footnotesize
\begin{center}
\begin{tabular}{@{}p{3.0cm}p{2.3cm}p{3.4cm}p{3.4cm}@{}}
    \toprule
    \textbf{Golongan} & \textbf{Register} & \textbf{Tanggung jawab} &
    \textbf{Akibat bila dilanggar} \\
    \midrule
    Boleh dirusak (\eng{caller-saved}) & \reg{AX}, \reg{BX},
    \reg{CX}, \reg{DX}, \reg{ES} & pemanggil menyimpan sendiri bila
    nilainya masih diperlukan & pemanggil memakai nilai yang sudah
    berubah \\
    \addlinespace
    Wajib dipulihkan (\eng{callee-saved}) & \reg{BP}, \reg{SI},
    \reg{DI}, \reg{DS}, \reg{SS} & prosedur menyimpan pada
    \eng{prologue} dan memulihkan pada \eng{epilogue} & pemanggil
    kehilangan nilai lama; akses data dapat salah alamat \\
    \addlinespace
    Selalu dipulihkan & \reg{SP} & stack harus seimbang saat
    \asm{RET} dijalankan & \asm{RET} mengambil alamat yang salah;
    program berhenti tak terkendali \\
    \bottomrule
\end{tabular}
\end{center}

\vspace{0.3em}
{\footnotesize\color{obeGray}Pada 8086 tidak ada aturan perangkat
keras yang memaksa, sehingga kesepakatan dibuat di tingkat program;
pembagian di atas mengikuti konvensi Microsoft C untuk 8086.\par}
\Sumber{Bab XIII, Konvensi Pemanggilan}
\end{frame}

% ---------------------------------------------------------
\begin{frame}{Nilai Balik Prosedur}
\small
Nilai balik dikembalikan melalui register, bukan melalui stack.

\vspace{0.3em}
\begin{center}
\begin{tabular}{@{}p{3.2cm}p{4.2cm}p{5.2cm}@{}}
    \toprule
    \textbf{Ukuran hasil} & \textbf{Register} & \textbf{Keterangan} \\
    \midrule
    \dec{8} bit & \reg{AL} &
    hasil satu byte, misalnya pencarian karakter \\
    \addlinespace
    \dec{16} bit & \reg{AX} &
    ukuran lazim untuk prosedur 8086 \\
    \addlinespace
    \dec{32} bit & \reg{DX}:\reg{AX} &
    hasil \asm{MUL} dan \asm{IMUL} memang tersimpan di kedua register
    ini \\
    \bottomrule
\end{tabular}
\end{center}

\vspace{0.35em}
\begin{block}{Mengapa bukan melalui stack}
\small \reg{DX}:\reg{AX} muncul secara alami karena \asm{MUL} dan
\asm{IMUL} menyimpan hasil kalinya pada kedua register itu.
\reg{AX} termasuk register yang boleh dirusak, sehingga prosedur
tidak wajib memulihkannya.
\end{block}
\Sumber{Bab XIII, Nilai Balik Prosedur}
\end{frame}

% ---------------------------------------------------------
\begin{frame}[fragile]{Rekursi pada Assembly}
\begin{columns}[t]
\begin{column}{0.55\textwidth}
\begin{kodeasm}[numbers=none]
Faktorial PROC
    PUSH BP              ; prologue
    MOV  BP, SP
    MOV  AX, [BP+4]      ; AX <- n
    CMP  AX, 1
    JBE  basis           ; kasus basis
    DEC  AX
    PUSH AX              ; dorong (n-1)
    CALL Faktorial       ; rekursi
    ADD  SP, 2           ; buang argumen
    MUL  WORD PTR [BP+4] ; AX <- (n-1)! * n
    POP  BP              ; epilogue
    RET
basis:
    MOV  AX, 1           ; 0! = 1! = 1
    POP  BP
    RET
Faktorial ENDP
\end{kodeasm}
\end{column}
\begin{column}{0.42\textwidth}
{\footnotesize
\begin{itemize}
    \item \term{Rekursi} adalah pemanggilan prosedur oleh dirinya
          sendiri.

    \vspace{0.2em}
    \item Pada bahasa tingkat tinggi rekursi tampak sebagai fitur
          bahasa; pada assembly rekursi tampak jelas sebagai
          penumpukan bingkai stack.

    \vspace{0.2em}
    \item Setiap pemanggilan menambah satu bingkai baru berisi
          parameter, alamat kembali, dan \reg{BP} sendiri.

    \vspace{0.2em}
    \item Tanpa kasus basis, rekursi tidak akan berhenti dan stack
          akan melampaui batasnya.
\end{itemize}}
\end{column}
\end{columns}
\end{frame}

% ---------------------------------------------------------
\begin{frame}{Kedalaman Rekursi dan \eng{Stack Overflow}}
\small
\begin{columns}[t]
\begin{column}{0.5\textwidth}
\begin{exampleblock}{Pemakaian stack}
\footnotesize
\begin{itemize}
    \item Setiap bingkai rekursif memerlukan \dec{6} byte:
          \dec{2} parameter, \dec{2} alamat kembali, \dec{2}
          \reg{BP} lama.
    \item Untuk $5!$ terbentuk lima bingkai $\approx$ \dec{30}
          byte, sehingga kebutuhan stack setara $O(n)$.
    \item Versi iteratif cukup satu bingkai saja ($O(1)$), tanpa
          biaya \asm{CALL} dan \eng{prologue} berulang.
\end{itemize}
\end{exampleblock}
\end{column}
\begin{column}{0.46\textwidth}
\begin{alertblock}{\eng{Stack overflow}}
\footnotesize Prosesor 8086 tidak memiliki pengaman batas stack. Bila
rekursi terlalu dalam atau kasus basis tidak pernah tercapai,
\reg{SP} terus turun melewati batas bawah ruang \dir{.STACK}, lalu
menimpa segmen data atau segmen kode --- tanpa pesan kesalahan.
\end{alertblock}
\end{column}
\end{columns}

\vspace{0.3em}
{\footnotesize\color{obeGray}Fibonacci rekursif $F(n) = F(n-1) +
F(n-2)$ jauh lebih berat: untuk $n = \dec{20}$ jumlah pemanggilannya
melebihi \dec{10.000}, sedangkan versi iteratifnya cukup satu
perulangan. Karena itu Fibonacci jauh lebih baik ditulis secara
iteratif pada assembly. Rekursi baru dipakai bila persoalannya memang
bercabang, misalnya penelusuran pohon.\par}
\Sumber{Bab XIII, Rekursi pada Assembly}
\end{frame}

% ---------------------------------------------------------
\begin{frame}{Kesalahan Umum Pengelolaan Stack}
\small Empat kesalahan berikut paling sering dijumpai.

\vspace{0.3em}
\begin{enumerate}
    \item \textbf{Stack tidak seimbang.} Jumlah \asm{PUSH} tidak sama
          dengan jumlah \asm{POP} pada setiap jalur eksekusi,
          sehingga \asm{RET} mengambil alamat yang salah.

    \vspace{0.25em}
    \item \textbf{\asm{RET} ke alamat yang salah.} Pemanggilan
          \eng{near} ditutup dengan \asm{RETF} atau sebaliknya,
          sehingga jumlah byte yang diambil dari stack tidak cocok.

    \vspace{0.25em}
    \item \textbf{Lupa menyimpan register yang dipakai.} Prosedur
          memakai \reg{SI} atau \reg{DI} tanpa menyimpannya, sehingga
          pemanggil kehilangan nilai lama register itu.

    \vspace{0.25em}
    \item \textbf{Kerusakan \eng{stack frame}.} \reg{BP} diubah di
          tengah prosedur, atau jalur keluar prosedur melompati
          \asm{POP BP}, sehingga \eng{epilogue} memulihkan nilai yang
          salah.
\end{enumerate}
\Sumber{Bab XIII, Kesalahan Umum Pengelolaan Stack}
\end{frame}

% ---------------------------------------------------------
\begin{frame}{Rangkuman Bagian 3}
\small
\begin{itemize}
    \item \eng{Stack frame} standar dibangun \eng{prologue}
          \asm{PUSH BP} dan \asm{MOV BP, SP}, lalu dibongkar
          \eng{epilogue} \asm{POP BP} dan \asm{RET}.

    \vspace{0.25em}
    \item Parameter diakses melalui \code{[BP+4]}, \code{[BP+6]}, dan
          seterusnya; variabel lokal melalui \code{[BP-2]} setelah
          \asm{SUB SP, n} pada \eng{prologue}.

    \vspace{0.25em}
    \item Konvensi pemanggilan mengatur letak parameter, letak nilai
          balik, dan tanggung jawab pembersihan stack --- melalui
          \asm{RET} \code{n} oleh prosedur, atau \asm{ADD SP, n} oleh
          pemanggil, tetapi tidak keduanya.

    \vspace{0.25em}
    \item Nilai balik diletakkan pada \reg{AL}, \reg{AX}, atau
          \reg{DX}:\reg{AX}. \reg{AX}, \reg{BX}, \reg{CX}, \reg{DX},
          dan \reg{ES} boleh dirusak; \reg{BP}, \reg{SI}, \reg{DI},
          \reg{DS}, dan \reg{SS} wajib dipulihkan.

    \vspace{0.25em}
    \item Rekursi pada assembly adalah penumpukan bingkai stack
          ($O(n)$) dan berisiko \eng{stack overflow}; versi iteratif
          memakai stack tetap ($O(1)$).
\end{itemize}
\end{frame}

% =========================================================
% CONTOH PROGRAM
% =========================================================
\begin{frame}[fragile]{Contoh 1: Prosedur \code{HitungLuas}}
\scriptsize Menghitung luas persegi panjang dengan prosedur
berparameter melalui stack; nilai balik pada \reg{AX}, prosedur
membersihkan sendiri argumennya dengan \asm{RET} \code{4}.

\vspace{0.15em}
\begin{columns}[t]
\begin{column}{0.53\textwidth}
\begin{kodeasm}[numbers=none]
.MODEL SMALL
.STACK 100h
.DATA
    lebar  DW 5
    tinggi DW 10
    luas   DW 0
.CODE
START:
    MOV  AX, @DATA
    MOV  DS, AX
    PUSH tinggi        ; param kedua
    PUSH lebar         ; param pertama
    CALL HitungLuas    ; AX <- luas
    MOV  luas, AX
    MOV  AH, 4Ch
    INT  21h
\end{kodeasm}
\end{column}
\begin{column}{0.44\textwidth}
\begin{kodeasm}[numbers=none]
HitungLuas PROC
    PUSH BP
    MOV  BP, SP
    MOV  AX, [BP+4]    ; lebar
    MUL  WORD PTR [BP+6]
    POP  BP
    RET  4             ; buang 4 byte
HitungLuas ENDP
END START
\end{kodeasm}

\vspace{0.2em}
\begin{block}{Hasil}
\footnotesize $5 \times 10 = \dec{50}$ (\hex{0032}) pada \reg{AX} dan
variabel \code{luas}. Pemanggil \emph{tidak} perlu \asm{ADD SP, 4}
karena \asm{RET} \code{4} sudah membuang kedua parameter.
\end{block}
\end{column}
\end{columns}

{\scriptsize\color{obeGray}Sumber: \file{prosedur-luas.asm} ---
prosedur ini hanya merusak \reg{AX}.\par}
\end{frame}

% ---------------------------------------------------------
\begin{frame}[fragile]{Contoh 2: Faktorial Rekursif}
\scriptsize Menghitung $5!$ secara rekursif; rekursi pada assembly
adalah penumpukan bingkai stack.

\vspace{0.1em}
\begin{columns}[t]
\begin{column}{0.4\textwidth}
\begin{kodeasm}[numbers=none, basicstyle=\ttfamily\fontsize{7.5}{8.4}\selectfont]
.MODEL SMALL
.STACK 100h
.DATA
    n     DW 5
    hasil DW 0
.CODE
START:
    MOV  AX, @DATA
    MOV  DS, AX
    PUSH n             ; argumen n
    CALL Faktorial     ; AX <- n!
    ADD  SP, 2
    MOV  hasil, AX
    MOV  AH, 4Ch
    INT  21h
\end{kodeasm}
\end{column}
\begin{column}{0.56\textwidth}
\begin{kodeasm}[numbers=none, basicstyle=\ttfamily\fontsize{7.5}{8.4}\selectfont]
Faktorial PROC
    PUSH BP              ; prologue
    MOV  BP, SP
    MOV  AX, [BP+4]      ; AX <- n
    CMP  AX, 1           ; n <= 1 ?
    JBE  basis           ; kasus basis
    DEC  AX
    PUSH AX              ; dorong (n-1)
    CALL Faktorial       ; rekursi
    ADD  SP, 2           ; buang argumen
    MUL  WORD PTR [BP+4] ; AX <- (n-1)! * n
    POP  BP              ; epilogue
    RET
basis:
    MOV  AX, 1           ; 0! = 1! = 1
    POP  BP
    RET
Faktorial ENDP
\end{kodeasm}
\end{column}
\end{columns}

{\scriptsize\color{obeGray}Hasil: \reg{AX} $=$ \dec{120} (\hex{0078});
lima bingkai $\times$ \dec{6} byte $\approx$ \dec{30} byte stack.
Sumber: \file{faktorial.asm}.\par}
\end{frame}

% =========================================================
% PRAKTIKUM
% =========================================================
\begin{frame}{Praktikum: Menulis Prosedur dan Memeriksa Stack}
\begin{block}{Judul praktikum}
\small Menulis prosedur berparameter dan memeriksa isi stack dengan
\eng{debugger} \quad\textbullet\quad \textbf{Sub-CPMK:} 3.4
\quad\textbullet\quad \textbf{Waktu:} 100 menit
\quad\textbullet\quad \textbf{Bentuk:} individu
\quad\textbullet\quad \textbf{Tugas 6}
\end{block}

\vspace{0.15em}
\begin{columns}[t]
\begin{column}{0.48\textwidth}
{\color{obeBlue}\bfseries Alat dan bahan}
\begin{itemize}\footnotesize
    \item Emu8086 dengan jendela register, memori, dan stack, atau
          GUI Turbo Assembler
    \item DOSBox beserta \file{DEBUG}, TASM, dan TLINK
    \item Berkas \file{prosedur-luas.asm} dan \file{faktorial.asm}
\end{itemize}
\end{column}
\begin{column}{0.48\textwidth}
{\color{obeBlue}\bfseries Langkah kerja}
\begin{enumerate}\footnotesize
    \item Rakit dan jalankan \file{prosedur-luas.asm}; pastikan
          \reg{AX} dan \code{luas} bernilai \hex{0032}.
    \item Ubah \code{lebar} menjadi \dec{7} dan \code{tinggi}
          \dec{6}; pastikan hasilnya \dec{42} (\hex{002A}).
    \item Telusuri isi stack per langkah; buktikan \code{[BP+4]}
          $=$ \hex{0005} dan \code{[BP+6]} $=$ \hex{000A}.
    \item Ubah \asm{RET} \code{4} menjadi \asm{RET} biasa pada salinan
          berkas, lalu jelaskan akibatnya.
\end{enumerate}
\end{column}
\end{columns}

\vspace{0.15em}
{\footnotesize\textbf{Dikumpulkan:} \file{prosedur-luas.asm}
berparameter tiga (memakai \code{[BP+8]}), \file{prosedur-ganjil.asm},
tabel pengamatan \reg{SP}/\reg{BP}/\reg{SS}, tangkapan layar jendela
stack, dan jawaban pertanyaan langkah 1(c), 3(b), 4(b).\par}
\end{frame}

% =========================================================
% LATIHAN
% =========================================================
\begin{frame}{Latihan}
\small
\begin{enumerate}
    \item Sebutkan tiga alasan mengapa program sebaiknya dipecah
          menjadi prosedur. Jelaskan pula mengapa pemakaian prosedur
          tidak selalu membuat program lebih cepat, dan kapan
          prosedur sebaiknya dihindari.

    \vspace{0.28em}
    \item Jelaskan apa yang dikerjakan \asm{CALL} dan \asm{RET}.
          Mengapa \asm{JMP} tidak dapat dipakai sebagai pengganti
          \asm{CALL}?

    \vspace{0.28em}
    \item Diketahui \reg{SS} $=$ \hex{2000} dan \reg{SP} $=$
          \hex{FFFC}. Tentukan alamat fisik sel teratas, lalu nilai
          \reg{SP} serta isi sel setelah \asm{MOV AX, 1234h},
          \asm{PUSH AX}, \asm{PUSH BX}, dan \asm{POP CX}. Berapa byte
          data yang masih tertinggal di stack?

    \vspace{0.28em}
    \item Prosedur \code{Jumlah} memakai \eng{prologue} \asm{PUSH BP}
          dan \asm{MOV BP, SP}, membaca \code{[BP+4]} dan
          \code{[BP+6]}, lalu ditutup \asm{RET} \code{4}. Tuliskan
          urutan instruksi pemanggil yang tepat, dan jelaskan mengapa
          pemanggil tidak boleh menambahkan \asm{ADD SP, 4}.

    \vspace{0.28em}
    \item Tulislah prosedur \code{HitungRataRata} yang menerima dua
          parameter 16-bit melalui stack dan mengembalikan hasil bagi
          keduanya (dibulatkan ke bawah) pada \reg{AX}, lengkap dengan
          \eng{prologue}, \eng{epilogue}, dan program pemanggilnya.
\end{enumerate}
\end{frame}

% =========================================================
% RANGKUMAN
% =========================================================
\begin{frame}{Rangkuman}
\footnotesize
\begin{itemize}
    \item Prosedur dipakai untuk memperoleh modularitas, penggunaan
          ulang, dan abstraksi. Biaya pemanggilannya berupa dorongan
          parameter, \asm{CALL}, \eng{prologue}, \eng{epilogue}, dan
          \asm{RET}.

    \vspace{0.18em}
    \item \asm{CALL} menyimpan alamat kembali \reg{CS}:\reg{IP} ke
          stack lalu melompat; \asm{RET} mengambil alamat itu dan
          mengembalikan kendali. Pemanggilan \eng{near} menyimpan
          \reg{IP} (\dec{2} byte) dan ditutup \asm{RET}, pemanggilan
          \eng{far} menyimpan \reg{CS} dan \reg{IP} (\dec{4} byte)
          dan ditutup \asm{RETF}.

    \vspace{0.18em}
    \item Stack bekerja secara LIFO dan ditunjuk \reg{SS}:\reg{SP},
          serta tumbuh ke arah alamat rendah: \asm{PUSH} menurunkan
          \reg{SP} \dec{2} byte, \asm{POP} menaikkannya kembali.
          \asm{PUSHF} dan \asm{POPF} menyimpan serta memulihkan
          \reg{FLAGS}.

    \vspace{0.18em}
    \item \eng{Stack frame} standar dibangun \eng{prologue}
          \asm{PUSH BP} dan \asm{MOV BP, SP}, lalu dibongkar
          \eng{epilogue} \asm{POP BP} dan \asm{RET}. Parameter diakses
          melalui \code{[BP+4]}, \code{[BP+6]}, dan seterusnya;
          variabel lokal melalui \code{[BP-2]}.

    \vspace{0.18em}
    \item Konvensi pemanggilan mengatur letak parameter, letak nilai
          balik, dan tanggung jawab pembersihan stack. Nilai balik
          diletakkan pada \reg{AL}, \reg{AX}, atau \reg{DX}:\reg{AX};
          \reg{AX}--\reg{DX} dan \reg{ES} boleh dirusak, sedangkan
          \reg{BP}, \reg{SI}, \reg{DI}, \reg{DS}, dan \reg{SS} wajib
          dipulihkan.

    \vspace{0.18em}
    \item Rekursi pada assembly adalah penumpukan bingkai stack.
          Rekursi memakai stack $O(n)$ dan berisiko
          \eng{stack overflow}, sedangkan versi iteratif biasanya
          memakai stack tetap $O(1)$.
\end{itemize}
\end{frame}

% ---------------------------------------------------------
\begin{frame}{Cek Pemahaman}
\small
\begin{enumerate}
    \item Instruksi \asm{CALL} \eng{near} menyimpan \dots\ ke dalam
          stack.
          \begin{enumerate}\footnotesize
              \item register \reg{CS} dan \reg{IP} \quad
              \item register \reg{IP} saja \quad
              \item seluruh register umum \quad
              \item register \reg{FLAGS} dan \reg{BP}
          \end{enumerate}

    \vspace{0.25em}
    \item Pada 8086, instruksi \asm{PUSH} bekerja dengan cara \dots
          \begin{enumerate}\footnotesize
              \item menambah \reg{SP} \dec{2}, lalu menyimpan data \quad
              \item mengurangi \reg{SP} \dec{2}, lalu menyimpan data \\
              \item menyimpan data, lalu menambah \reg{SP} \dec{2} \quad
              \item menyimpan data, lalu mengurangi \reg{SP} \dec{2}
          \end{enumerate}

    \vspace{0.25em}
    \item Setelah \eng{prologue} dijalankan, parameter pertama yang
          didorong paling akhir oleh pemanggil berada pada \dots
          \begin{enumerate}\footnotesize
              \item \code{[BP+0]} \quad \item \code{[BP+2]} \quad
              \item \code{[BP+4]} \quad \item \code{[BP+6]}
          \end{enumerate}
\end{enumerate}
\end{frame}

% =========================================================
% PENUTUP
% =========================================================
\begin{frame}{Penutup}
\small
\begin{block}{Sebelum pertemuan berikutnya}
\begin{itemize}
    \item Selesaikan Tugas 6: \file{prosedur-luas.asm} berparameter
          tiga, \file{prosedur-ganjil.asm}, tabel pengamatan
          \reg{SP}/\reg{BP}/\reg{SS}, dan tangkapan layar jendela
          stack.
    \item Ulangi bagian yang belum dikuasai dengan menelusuri
          \file{prosedur-luas.asm} dan \file{faktorial.asm} satu
          instruksi demi satu pada \eng{debugger}.
    \item Baca \textbf{Bab XIV} --- Interupsi BIOS/DOS dan Operasi
          I/O.
\end{itemize}
\end{block}

\vspace{0.35em}
{\footnotesize\textbf{Bacaan lanjutan:} Bab XIII bagian Konvensi
Pemanggilan dan Rekursi pada Assembly; pembahasan mode pengalamatan
\eng{base-relative} pada Bab VIII; pembahasan segmen pada Bab VI.\par}

\vspace{0.4em}
\begin{center}
{\color{obeBlue}\sffamily\large\bfseries Terima kasih\par}
\vspace{0.12em}
{\color{obeGray}\sffamily\small \NamaMK\ (\KodeMK) \textbullet\
\SemesterTeks\par}
\end{center}
\end{frame}

\end{document}
